\documentclass[11pt]{article}
\usepackage[a4paper,margin=27mm]{geometry}
\usepackage{amsmath,amssymb,amsthm,mathtools}
\usepackage[T1]{fontenc}
\usepackage{lmodern}
\usepackage[hidelinks]{hyperref}
\usepackage{microtype}
\newtheorem{theorem}{Theorem}[section]
\newtheorem{proposition}[theorem]{Proposition}
\newtheorem{lemma}[theorem]{Lemma}
\theoremstyle{remark}\newtheorem{remark}[theorem]{Remark}
\newcommand{\E}{\mathbb E}
\newcommand{\Pp}{\mathbb P}
\newcommand{\R}{\mathbb R}
\newcommand{\TV}{\mathrm{TV}}
\newcommand{\one}{\mathbf 1}
\newcommand{\sign}{\operatorname{sign}}
\newcommand{\Cov}{\operatorname{Cov}}
\newcommand{\diag}{\operatorname{diag}}
\newcommand{\cN}{\mathcal N}
\newcommand{\norm}[1]{\left\lVert #1\right\rVert}
\title{An explicit superpolynomial upper bound for the majority-consensus threshold}
\author{A proof note based on Kanoria--Montanari (2011)}
\date{September 7, 2026}
\begin{document}
\maketitle
\begin{abstract}
For the synchronous majority dynamics and fair tie-breaking convention of
Kanoria--Montanari, this note derives the upper bound
\[
 \theta_*(k)\le
 \exp\!\left\{-\frac{\log k}{\log 2}
  \log\log k+O(\log k)\right\}.
\]
The argument makes the time horizon in the dynamic cavity calculation
quantitative. Its two main ingredients are bounds on the covariance and
response of the cavity process, and a one-coordinate local central limit
estimate that does not involve the least probable trajectory.
This is an upper bound, not an asymptotic equivalence; no matching
large-degree lower bound is asserted.
\end{abstract}

\section{Statement and relationship to the original paper}
All logarithms are natural. The degree is denoted by $k$; the initial spin
mean is $\theta$, so that $\Pp(\sigma_v(0)=+1)=(1+\theta)/2$.
The update is synchronous majority, with a fresh fair output coin at a tie.
This has the same law as the tie-breaking convention in \cite{KM}.
We retain the paper's definitions of $C(t,s)$, $R(t,s)$ and the unbiased
cavity process $\tau$:
\[
 \tau(t+1)=\sign\!\left(\eta(t)+\sum_{s<t}R(t,s)\tau(s)\right),
 \qquad \E\eta(t)\eta(s)=C(t,s),
\]
where $\tau(0)$ is a fair sign independent of the Gaussian process $\eta$.
The consistency conditions are
\[
 C(t,s)=\E\tau(t)\tau(s),\qquad
 R(t,s)=\left.\partial_{h(s)}\E\tau_h(t)\right|_{h=0}.
\]
Put $C_T=(C(t,s))_{0\le t,s\le T}$,
$R_T=(R(t,s))_{0\le t,s\le T}$, with zero entries for $s\ge t$, and
\[
 r_j=R(j,j-1),\qquad a_0=1,\qquad a_t=\prod_{j=1}^t r_j.
\]
The existence, strict positive definiteness, and finite-time cavity
identities proved in \cite{KM} are used below. The proof does not substitute
a growing time into an unquantified fixed-time convergence theorem.

\begin{theorem}\label{thm:finite}
There are absolute constants $A,B,k_0$ such that the following holds.
Let $T\ge0$ and $k\ge k_0$ satisfy
\begin{equation}\label{eq:degree-condition}
 \log k\ge A2^{T/2}.
\end{equation}
For
\begin{equation}\label{eq:initial-bias}
 \theta=\frac{8}{a_T k^{(T+1)/2}},
\end{equation}
one has $0<\theta<1$ and
\[
 \Pp_\theta\{\sigma_v(T+2)=-1\}\le e^{-k/8}.
\]
Consequently,
\begin{equation}\label{eq:finite-bound}
 \theta_*(k)\le\frac{8}{a_T k^{(T+1)/2}}
 \le k^{-(T+1)/2}\exp\{B2^{T/2}\}.
\end{equation}
The conclusion holds for odd and even degrees.
\end{theorem}

\begin{theorem}\label{thm:asymptotic}
As $k\to\infty$,
\begin{equation}\label{eq:main}
 \theta_*(k)\le
 \exp\!\left\{-\frac{\log k}{\log 2}
  \log\log k+O(\log k)\right\}.
\end{equation}
In particular, for every $c<1/\log 2$ and every sufficiently large $k$,
\[
 \theta_*(k)\le e^{-c\log k\log\log k}.
\]
\end{theorem}

The original Theorem 2.3 in \cite{KM} gives $\theta_*(k)=o(k^{-M})$
for every fixed $M$. The improvement in \eqref{eq:main} is an explicit
choice of a diverging exponent. The constant $1/\log 2$ is a constant
produced by this argument; it is not claimed to be the true asymptotic
constant of the threshold.

\section{Structural facts about the cavity process}
\begin{lemma}\label{lem:structure}
Every finite cavity trajectory is associated and invariant under
simultaneous spin reversal. Moreover,
\[
 C(t,s)\ge0,\qquad R(t,s)\ge0,
\]
and
\begin{equation}\label{eq:parity}
 C(t,s)=0\quad(t-s\text{ odd}),\qquad
 R(t,s)=0\quad(t-s\text{ even}).
\end{equation}
The even and odd spin subsequences are independent. For every nonempty
finite set $J$ of time indices,
\begin{equation}\label{eq:constant-path-mass}
 \Pp\{\tau(s)=+1\text{ for all }s\in J\}\ge2^{-|J|},\qquad
 C_J\succeq 2^{1-|J|}\one_J\one_J^\top.
\end{equation}
\end{lemma}
\begin{proof}
Induct on time. If the previously constructed responses are nonnegative,
the effective-process spins are nondecreasing functions of the past
Gaussian fields, the initial sign, and the added fields $h$.
Thus the newly constructed responses are nonnegative. A Gaussian vector
with nonnegative covariances is associated \cite{Pitt}; adjoining an
independent initial sign preserves association. Applying nondecreasing
maps therefore shows that the enlarged spin trajectory is associated,
and hence that its covariances are nonnegative. Simultaneous spin reversal
preserves its law and gives zero means.

The parity identities follow from the recursion. The odd spins depend
only on the even-indexed Gaussian fields. The even spins depend only on
$\tau(0)$ and the odd-indexed Gaussian fields. The two Gaussian blocks are
independent, giving the asserted independence and completing the parity
induction. This is the bipartite decomposition noted in \cite{KM}.

Association gives the first inequality in \eqref{eq:constant-path-mass}.
The two constant sign vectors have equal probabilities. Their contribution
to $\E\tau_J\tau_J^\top$ gives the second inequality.
\end{proof}

\begin{lemma}[Gaussian projection inequality]\label{lem:projection}
Let $\boldsymbol c_T=(C(t,0))_{0\le t\le T}$. Then
\begin{equation}\label{eq:projection}
 C_T\succeq \boldsymbol c_T\boldsymbol c_T^\top+R_TC_TR_T^\top.
\end{equation}
In particular,
\begin{equation}\label{eq:memory-bounds}
 \E\left(\sum_{s<t}R(t,s)\tau(s)\right)^2\le1,
 \qquad \sum_{s<t}R(t,s)^2\le1,
 \qquad \sum_{s<t}R(t,s)\le\sqrt t.
\end{equation}
\end{lemma}
\begin{proof}
Shifting the Gaussian vector by $h$ is exactly the same as adding $h$ in
the effective-process recursion. Gaussian differentiation therefore gives
\[
 \E[\tau^T(\eta^T)^\top]=R_TC_T.
\]
The orthogonal projection of $\tau^T$ onto the linear span of
$\tau(0),\eta(0),\ldots,\eta(T)$ is
$\boldsymbol c_T\tau(0)+R_T\eta^T$. The covariance matrix of its residual is
positive semidefinite, proving \eqref{eq:projection}.
Its diagonal gives the first inequality in \eqref{eq:memory-bounds}.
Since $R$ and $C$ are entrywise nonnegative and $C(s,s)=1$, it also gives
$\sum_sR(t,s)^2\le1$. Cauchy--Schwarz proves the final assertion.
\end{proof}

\section{A Gaussian boundary-mass estimate}
The following lemma is the device that avoids a minimum-trajectory
probability. A rare, arbitrary sign pattern is not used.

\begin{lemma}\label{lem:boundary}
Let $X\in\{-1,+1\}^m$, let $C=\E XX^\top$ be positive definite, and
suppose
\[
 p=\Pp\{X_1=\cdots=X_m\}>0.
\]
For $G\sim\cN(0,C)$ and $m\ge2$,
\begin{equation}\label{eq:conditional-orthant}
 \Pp\{G_1,\ldots,G_{m-1}\ge0\mid G_m=0\}
 \ge \frac{e^{-1}p^{(m-2)/2}}
 {2^{3(m-1)/2}\Gamma((m+1)/2)}.
\end{equation}
For $m=1$ the conditional orthant probability is understood as $1$.
Furthermore, writing $\phi_{\mu,C}$ for Gaussian density and
$L=\mu^\top C^{-1}\mu$, one has
\begin{equation}\label{eq:shifted-face}
 \int_{\R_+^{m-1}}\phi_{\mu,C}(z_1,\ldots,z_{m-1},0)\,dz
 \ge (2\pi)^{-1/2}e^{-L}
 \left[\Pp\{G_{<m}\ge0\mid G_m=0\}\right]^2.
\end{equation}
\end{lemma}
\begin{proof}
Set $D_i=(1-X_iX_m)/2$ for $i<m$, $q=\E D$, and $M=\E DD^\top$.
A direct calculation gives
\[
 C_{<m,<m}-C_{<m,m}C_{m,<m}=4(M-qq^\top).
\]
This matrix is positive definite. Consider jointly centered Gaussian
$(W,Z)$ with covariance
\[
 \begin{pmatrix} M&q\\q^\top&1\end{pmatrix}.
\]
This covariance matrix is the Gram matrix of the random vector $(D,1)$.
Let $v=1-q^\top M^{-1}q$. For every deterministic vector $a$,
\[
 \E(1-a^\top D)^2\ge\Pp(D=0)=p,
\]
so that $v\ge p$. Conditional on $W$,
$Z$ has mean $h(W)=q^\top M^{-1}W$ and variance $v$.
The variance of $h(W)$ is $1-v\le1$.

Write $W=M^{1/2}Z_0$, where $Z_0$ is a standard Gaussian vector
in dimension $d=m-1$. Then $h(W)=a^\top Z_0$ for a vector with
$\norm a^2=1-v\le1$. The entries of $M$ are nonnegative, so Gaussian
association gives $\Pp(W\ge0)\ge2^{-d}$. The event $W\ge0$ is a cone
in $Z_0$, and hence is independent of the Gaussian radius $\norm{Z_0}$.
On $\{\norm{Z_0}\le\sqrt v\}$, the density ratio for conditioning $Z=0$
satisfies
\[
 \frac{f_{Z\mid W}(0\mid W)}{f_Z(0)}
 =v^{-1/2}\exp\{-h(W)^2/(2v)\}\ge v^{-1/2}e^{-1/2}.
\]
Since $0<v\le1$, integration of the standard Gaussian density over the
ball of radius $\sqrt v$ gives
\[
 \Pp(\norm{Z_0}\le\sqrt v)
 \ge \frac{e^{-1/2}v^{d/2}}{2^{d/2}\Gamma(d/2+1)}.
\]
Combining these facts and using $v\ge p$ yields
\[
 \Pp(W\ge0\mid Z=0)
 \ge \frac{e^{-1}v^{(d-1)/2}}{2^{3d/2}\Gamma(d/2+1)}
 \ge \frac{e^{-1}p^{(m-2)/2}}
             {2^{3(m-1)/2}\Gamma((m+1)/2)}.
\]
The conditional covariance of $W$ is $M-qq^\top$; multiplication by $2$
identifies its orthant probability with the one in
\eqref{eq:conditional-orthant}.

For the second assertion, put
$\Gamma=C_{<m,<m}-C_{<m,m}C_{m,<m}$ and
$\nu=\mu_{<m}-C_{<m,m}\mu_m$.
The conditional mean and covariance on the last-coordinate-zero slice
are $\nu,\Gamma$, and
\[
 L=\mu_m^2+\nu^\top\Gamma^{-1}\nu.
\]
For any measurable $A$, Cauchy--Schwarz applied to the Gaussian
likelihood ratio gives
\[
 \cN(\nu,\Gamma)(A)
 \ge \cN(0,\Gamma)(A)^2
       e^{-\nu^\top\Gamma^{-1}\nu}.
\]
Multiplying by the marginal density
$(2\pi)^{-1/2}e^{-\mu_m^2/2}$ proves \eqref{eq:shifted-face}.
\end{proof}

\section{Quantitative bounds on response and covariance}
\begin{proposition}\label{prop:quantitative-cavity}
There is an absolute constant $B$ such that, for every $t\ge1$,
\begin{equation}\label{eq:response-lower}
 0<r_t\le1,\qquad r_t\ge\exp\{-B2^{t/2}\},
\end{equation}
and, for every $T\ge0$,
\begin{equation}\label{eq:covariance-lower}
 \lambda_{\min}(C_T)\ge\exp\{-B2^{T/2}\},\qquad
 a_T\ge\exp\{-B2^{T/2}\}.
\end{equation}
\end{proposition}
\begin{proof}
First consider $r_{t+1}=R(t+1,t)$. Let
\[
 I=\{s\le t:s\equiv t\pmod2\},\qquad
 J=\{s<t:s\not\equiv t\pmod2\}.
\]
Set $n=|I|=\lfloor t/2\rfloor+1$ and $m=|J|=\lceil t/2\rceil$.
To lower-bound the last-step response, retain only the trajectory in which
all the past spins in the active parity class are $+1$.
On that trajectory the Gaussian field vector on $I$ has deterministic
shift
\[
 \mu=R_{I,J}\one_J.
\]
If $J$ is empty, $\mu=0$. Otherwise, \eqref{eq:projection} and
\eqref{eq:constant-path-mass} imply
\[
 C_I\succeq R_{I,J}C_JR_{I,J}^\top
 \succeq 2^{1-m}\mu\mu^\top,
 \qquad \mu^\top C_I^{-1}\mu\le2^{m-1}.
\]
The spin vector $\tau_I$ satisfies
$\Pp(\tau_I\text{ is constant})\ge2^{1-n}$.
Lemma \ref{lem:boundary} gives a conditional Gaussian orthant
probability of at least $\exp\{-B(n+1)^2\}$. Its shifted-face bound and
$\mu^\top C_I^{-1}\mu\le2^{m-1}$ therefore lower-bound the Gaussian
face mass for this trajectory by
\[
 \exp\{-B(2^m+(n+1)^2)\}\ge\exp\{-B'2^{t/2}\}.
\]
Here the empty-$J$ case has zero shift cost and is included by enlarging
the absolute constants.
This face mass is a lower bound on $r_{t+1}$: when the active class
contains $\tau(0)$ its probability $1/2$ is canceled by the factor $2$
in differentiating the expectation of a sign; otherwise the factor is $2$
and may simply be discarded. This proves the lower bound in
\eqref{eq:response-lower}, after changing the absolute constant.
The upper bound follows from \eqref{eq:memory-bounds}.

For the covariance bound, iterate \eqref{eq:projection}, using the
nilpotence of $R_T$:
\[
 C_T\succeq\sum_{j=0}^T R_T^j \boldsymbol c_T\boldsymbol c_T^\top(R_T^\top)^j=V_TV_T^\top,
 \qquad V_T=(R_T^j\boldsymbol c_T)_{j=0}^T.
\]
The displayed vectors are the columns of $V_T$. This matrix is lower
triangular and its $j$th diagonal
entry, with indexing starting at $0$, is $a_j$. Consequently,
\[
 \det C_T\ge\prod_{j=0}^T a_j^2
 =\prod_{j=1}^T r_j^{2(T+1-j)}.
\]
Since $\lambda_{\max}(C_T)\le T+1$,
\[
 -\log\lambda_{\min}(C_T)
 \le T\log(T+1)+2\sum_{j=1}^T(T+1-j)(-\log r_j).
\]
The geometric-sum estimates
\[
 \sum_{j=1}^T2^{j/2}\le B2^{T/2},\qquad
 \sum_{j=1}^T(T+1-j)2^{j/2}\le B2^{T/2}
\]
prove both assertions in \eqref{eq:covariance-lower}.
\end{proof}

\begin{remark}
The factor $2^{T/2}$, rather than $2^T$, is important. The boundary
calculation involves one parity block. Discarding this decomposition
would lose a factor $2$ in the leading coefficient of \eqref{eq:main}.
\end{remark}

\section{A one-coordinate local central limit estimate}
The original local CLT assumes a lower bound on every trajectory atom.
Here only one coordinate is pinned, allowing that assumption to be
removed.

\begin{lemma}\label{lem:one-face-clt}
Let $X_1,\ldots,X_N$ be i.i.d. in $\{-1,+1\}^d$, with mean $m$ and
covariance $\Sigma$, where
\[
 \norm{\sqrt N m}_\infty\le M,\qquad \Sigma\succeq\lambda I,
 \qquad 0<\lambda\le1.
\]
Let $S=\sum_iX_i$, $\mu=\sqrt N m$ and $G\sim\cN(\mu,\Sigma)$.
There is a fixed numerical exponent $c_0$ such that, with
\[
 H=\left[\frac{C(d+1)(M+1)}\lambda\right]^{c_0},
\]
the following estimates hold whenever $N\ge H^2$.
Ordinary orthant probabilities, with integer thresholds of absolute value
at most $5$, have Gaussian approximation error at most $H/\sqrt N$,
after normalizing $S$ by $\sqrt N$.
For any coordinate $j$, any admissible parity value $b$ with $|b|\le5$,
and signs $\varepsilon_i\in\{-1,+1\}$,
\begin{align}
 &\left|\frac{\sqrt N}{2}\Pp\{S_j=b,
       \ \varepsilon_i S_i\ge b_i\ (i\ne j)\}
       -\int_{\substack{z_j=0\\\varepsilon_i z_i\ge0\ (i\ne j)}}
           \phi_{\mu,\Sigma}(z)\,dz_{-j}\right|
 \le\frac H{\sqrt N},\label{eq:face-clt}
\end{align}
where $|b_i|\le5$. The same statement holds with strict inequalities
and with averages of strict and weak inequalities.
Gaussian orthant probabilities and these Gaussian face integrals are
Lipschitz in their mean and covariance on the indicated parameter range,
with a constant bounded by a fixed polynomial in
$d+1,M+1,\lambda^{-1}$.
\end{lemma}
\begin{proof}
The ordinary approximation follows from the convex-set multivariate
Berry--Esseen bound \cite{Bentkus}. After whitening, each centered
summand has norm at most $2\sqrt{d/\lambda}$, giving an error at most
$C d^{7/4}\lambda^{-3/2}/\sqrt N$.

For \eqref{eq:face-clt}, condition on $S_j=b$. There are exactly
$N_+=(N+b)/2$ summands with $j$th coordinate $+1$ and
$N_-=(N-b)/2$ with $j$th coordinate $-1$.
Conditional on these counts, the sum of the remaining coordinates has
the distribution of a sum of $N_+$ independent vectors with law
$X_{-j}\mid X_j=+1$ and $N_-$ independent vectors with law
$X_{-j}\mid X_j=-1$.
It does not matter which summand indices have the two signs.

Write $p_\pm=\Pp(X_j=\pm1)$ and let $\Sigma_\pm$ be these two conditional
covariances. The weighted average
\[
 p_+\Sigma_++p_-\Sigma_-
 =\Sigma_{-j,-j}-\frac{\Sigma_{-j,j}\Sigma_{j,-j}}{\Sigma_{jj}}
 =:\Gamma
\]
is a Schur complement and satisfies $\Gamma\succeq\lambda I$.
Changing the weights from $p_\pm$ to $N_\pm/N$ changes this matrix
in operator norm by at most $C d(M+1)/\sqrt N$.
Thus the conditional sum covariance is at least $N\lambda I/2$ under
the stated size assumption. Apply the independent, not necessarily
identically distributed version of the same Berry--Esseen bound.

Since $X_j$ is binary, its conditional regression is exactly affine.
The conditional mean of the normalized sum is therefore
\[
 \mu_{-j}+\frac{\Sigma_{-j,j}}{\Sigma_{jj}}
                (b/\sqrt N-\mu_j),
\]
which is the corresponding Gaussian conditional mean.
The binomial local estimate, obtained directly from Stirling's formula,
is
\[
 \Pp(S_j=b)=\frac2{\sqrt N}
       \phi_{\mu_j,\Sigma_{jj}}(b/\sqrt N)+O(H/N).
\]
Multiplying this estimate by the conditional orthant approximation proves
\eqref{eq:face-clt}. Moving the bounded thresholds to zero has the same
order of error.

For completeness, covariance and mean comparison can be done on Gaussian
densities before integration. A mean derivative contributes a linear
Gaussian score and a covariance derivative contributes a quadratic score.
Their absolute integrals on an orthant are bounded by their full-space
absolute integrals. On a coordinate-zero face, condition on that
coordinate first. The conditional mean and covariance are the affine
regression and Schur complement above; their inverse norms are bounded
by a fixed polynomial in $\lambda^{-1}$, and their means by a fixed
polynomial in $M+1,\lambda^{-1}$. The first two conditional Gaussian
moments consequently give the asserted polynomial Lipschitz bounds.
All powers needed in this proof are fixed numbers, so one may enlarge
$c_0$ once to cover all estimates.
\end{proof}

\section{Kernel calculus without high-order product expansions}
For $\ell\in\{-1,+1\}$ define
\[
 \kappa_\ell(z)=\one\{\ell z>0\}+\tfrac12\one\{z=0\}.
\]
For $\lambda=(\lambda_0,\ldots,\lambda_d)$ and $u=(u_0,\ldots,u_{d-1})$,
put
\[
 F_{\lambda,u}(z)=\prod_{s=0}^{d-1}\kappa_{\lambda_{s+1}}(z_s+u_s).
\]
For a Gaussian mean $\mu$ and covariance $\Sigma$, write
\[
 J_s(\lambda;\mu,\Sigma)
 =\int_{\substack{z_s=0\\\lambda_{r+1}z_r\ge0\ (r\ne s)}}
             \phi_{\mu,\Sigma}(z)\,dz_{-s}.
\]
In the cavity process at dimension $d=t+1$, use the abbreviation
\[
 \mu_r(\lambda)=\sum_{v<r}R(r,v)\lambda_v,\qquad
 J_s^t(\lambda)=J_s(\lambda;\mu(\lambda),C_t).
\]
Gaussian differentiation of the path formula gives
\begin{equation}\label{eq:response-face-identity}
 R(t+1,s)=\frac12\sum_{\lambda\in\{-1,+1\}^{t+2}}
          \lambda_{t+1}\lambda_{s+1}J_s^t(\lambda).
\end{equation}
In particular $J_t^t$ is independent of $\lambda_{t+1}$ and
$\sum_{\lambda_0,\ldots,\lambda_t}J_t^t(\lambda)=r_{t+1}$.
Also $0\le J_s^t(\lambda)\le(2\pi)^{-1/2}$.

\begin{lemma}\label{lem:kernel-calculus}
Under the hypotheses of Lemma \ref{lem:one-face-clt}, changing a bounded
integer field $u$ to $v$ gives
\[
 \E[F_{\lambda,v}(S)-F_{\lambda,u}(S)]
 =\frac1{\sqrt N}\sum_s(v_s-u_s)\lambda_{s+1}
          J_s(\lambda;\mu,\Sigma)
  +O\!\left(\frac HN\norm{v-u}_1\right).
\]
Here all field coordinates have absolute value at most $2$.
For two stochastically ordered laws $P\preceq Q$ on $\{-1,+1\}^d$,
let $P_\alpha=(1-\alpha)P+\alpha Q$, and put
$\Delta_s=\E_QX_s-\E_PX_s\ge0$.
The exact interpolation identity is
\begin{align}
 &\E_{Q^{\otimes n}}F\left(\sum_iX_i\right)
       -\E_{P^{\otimes n}}F\left(\sum_iX_i\right)\nonumber\\
 &\quad=n\int_0^1
    \E\left[F(S_\alpha+Y)-F(S_\alpha+X)\right]d\alpha,
 \label{eq:product-interpolation}
\end{align}
where $(X,Y)$ is any monotone coupling of $P,Q$, independent of
$S_\alpha$, and $S_\alpha$ is the sum of $n-1$ i.i.d. $P_\alpha$ vectors.
Consequently its leading term is
\[
 \sqrt n\sum_s\Delta_s\lambda_{s+1}
        \int_0^1J_s(\lambda;\mu_\alpha,\Sigma_\alpha)\,d\alpha,
\]
with error bounded by $C H\sum_s\Delta_s$ when the means and covariances
throughout the interpolation satisfy the local CLT bounds.
Replacing $n/\sqrt{n-1}$ by $\sqrt n$ is absorbed in this error.
\end{lemma}
\begin{proof}
Telescope the product defining $F$ over its time coordinates.
At a changed coordinate the difference of the two majority kernels is
supported on a bounded number of lattice levels. Its signed total weight
on the lattice is $\lambda_{s+1}(v_s-u_s)/2$. Each level has the $2/\sqrt N$ prefactor in
\eqref{eq:face-clt}. At other coordinates the factors are strict or weak
half-line indicators or averages of these, with bounded threshold shifts.
Apply \eqref{eq:face-clt} to each term. This proves the first assertion;
it also covers a change of a child trajectory by taking fields $u+X,u+Y$.

Differentiate $P_\alpha^{\otimes n}$ and use permutation symmetry of the
summands to obtain \eqref{eq:product-interpolation}. Under the monotone
coupling, $\E\norm{Y-X}_1=\sum_s\Delta_s$.
Apply the first assertion conditionally on $X,Y$ and integrate.
No division by a path probability is made, and no expansion into terms
with two or more replaced children is needed.
\end{proof}

\section{Uniform finite-degree cavity estimates}
Let $P_{k,\theta}^{u}$ denote the root trajectory law in the rooted tree
with $k-1$ children and imposed parent trajectory $u$, as in \cite{KM}.
Only the displayed finite prefix of this law is used. Parent coordinates
may be $-1,0,+1$. For a prescribed root trajectory $\lambda$, the exact
cavity recursion is
\begin{equation}\label{eq:exact-cavity}
 P_{k,\theta}^{u}(\lambda)
 =\frac{1+\theta\lambda_0}{2}
   \E\left[F_{\lambda,u}\left(\sum_{i=1}^{k-1}X_i\right)\right],
\end{equation}
where the $X_i$ are independent with law $P_{k,\theta}^{\lambda}$ on the
appropriate shorter child trajectories. The notation denotes dynamics
under an imposed trajectory, not conditioning on an observed trajectory.

In the estimates below, a constant of the form
\begin{equation}\label{eq:controlled-constant}
 \exp\{C2^{T/2}\}
\end{equation}
will be called a controlled constant. Products of a fixed number of such
constants, multiplication by $2^{CT}(T+1)^C$, and sums over $t\le T$
remain of this form. Also, if
$L_t\le\exp\{C2^{t/2}\}$, then
\begin{equation}\label{eq:product-constant}
 \prod_{t=0}^T(2L_t)
 \le\exp\{C'2^{T/2}\}.
\end{equation}
This follows by summing a geometric series; it does not cost an extra
factor $T$ in the exponent.

\begin{proposition}\label{prop:unbiased-uniform}
There is an absolute $C$ such that, for all $T$ and sufficiently large
$k$ with $\log k\ge C2^{T/2}$, uniformly in $u$,
\begin{align}
 &\norm{P_{k,0}^{u}(\sigma^T\in\cdot)-\mathcal L(\tau^T)}_{\TV}
 \le \frac{\exp\{C2^{T/2}\}}{\sqrt k},\label{eq:unbiased-tv}\\
 &\max_{t\le T}\left|\sqrt k\,\E_{P_{k,0}^{u}}\sigma(t)
                 -\sum_{s<t}R(t,s)u_s\right|
 \le \frac{\exp\{C2^{T/2}\}}{\sqrt k}.
 \label{eq:unbiased-response}
\end{align}
\end{proposition}
\begin{proof}
Let $e_t$ be the supremum over $u$ of the sum of the two errors on the
left sides, at horizon $t$. Initially $e_0=0$.
For prescribed $\lambda$, the child means in \eqref{eq:exact-cavity}
obey
\[
 \sqrt{k-1}\,\E X_s=\mu_s(\lambda)+O(e_t+(t+1)/k).
\]
Their covariance differs from $C_t$ entrywise by
$O(e_t+(t+1)/k)$. This follows from the total-variation bound on raw
second moments and the $O(\sqrt{t+1}/\sqrt k)$ bound on the means.
Under a bootstrap $e_t\le\lambda_{\min}(C_t)/(100(t+1))$, the covariance
is at least $C_t/2$ in the needed eigenvalue sense and the normalized
means are bounded by $2(t+2)$.

Apply the ordinary orthant CLT in \eqref{eq:exact-cavity}, followed by
Gaussian parameter comparison. Summing over $2^{t+2}$ root trajectories
bounds the new total-variation error by $L_t(e_t+k^{-1/2})$, where
\[
 L_t\le 2^{C(t+1)}
      \left[\frac{C(t+2)}{\lambda_{\min}(C_t)}\right]^{c_1}
 \le\exp\{C'2^{t/2}\}
\]
for a fixed exponent $c_1$.

For the normalized mean, subtract the recursion with $u=0$ from that
with $u$. The zero-field mean is exactly zero by spin-flip symmetry.
The first part of Lemma \ref{lem:kernel-calculus}, multiplied by the
initial factor $1/2$, gives the leading term
\[
 \frac{1}{2\sqrt k}\sum_s u_s\lambda_{s+1}J_s^t(\lambda)
\]
for the trajectory-probability difference. Multiply by $\lambda_{t+1}$,
sum over trajectories, and use \eqref{eq:response-face-identity}. The result is
\[
 k^{-1/2}\sum_sR(t+1,s)u_s.
\]
The scaled error is at most $L_t(e_t+k^{-1/2})$.
Thus, after increasing $L_t$,
\[
 e_{t+1}\le L_t(e_t+k^{-1/2}).
\]
Equation \eqref{eq:product-constant} proves the claimed estimates.
The degree condition, with a sufficiently large absolute constant,
makes every obtained $e_t$ smaller than the bootstrap threshold and
satisfies the size requirements in the local CLT. This closes the
induction.
\end{proof}

\begin{proposition}[Uniform weak-bias amplification]\label{prop:bias-uniform}
Fix $T$ and let $\theta$ be as in \eqref{eq:initial-bias}. Under a degree
condition of the form \eqref{eq:degree-condition}, uniformly for $t\le T$
and imposed parent trajectories $u$,
\begin{equation}\label{eq:bias-amplification}
 d_t(u):=\E_{P_{k,\theta}^{u}}\sigma(t)
          -\E_{P_{k,0}^{u}}\sigma(t)
 =\theta a_t k^{t/2}(1+\varepsilon_t(u)),
 \qquad
 |\varepsilon_t(u)|\le
       \frac{\exp\{C2^{T/2}\}}{\sqrt k}.
\end{equation}
\end{proposition}
\begin{proof}
Write $b_t=\theta a_t k^{t/2}$ and
$\Lambda_T=\max(1,\max_{j\le T}r_j^{-1})$.
By Proposition \ref{prop:quantitative-cavity}, $\Lambda_T$ is a controlled
constant. Increasing the degree constant ensures
$\sqrt k\ge2\Lambda_T$. Hence $b_t$ is increasing, and
\begin{equation}\label{eq:bias-size-bounds}
 b_T=\frac8{\sqrt k},\qquad
 b_t\le\frac{8\Lambda_T}{k}\ (t<T),\qquad
 \sum_{s<t}b_s\le\frac{2\Lambda_T}{\sqrt k}b_t.
\end{equation}
The assertion at $t=0$ is exact. Bootstrap that
$|d_s(u)/b_s-1|\le1/2$ through time $t<T$.

For a prescribed root trajectory $\lambda$, apply
\eqref{eq:product-interpolation} to the child laws
$P=P_{k,0}^{\lambda}$ and $Q=P_{k,\theta}^{\lambda}$.
They admit a monotone coupling using the same initial uniforms and fair
output coins. Put $D_t=\max_{s\le t,u}d_s(u)$.
For the interpolated child laws, the normalized-mean difference from
$\mu(\lambda)$ is at most
\[
 C\bigl(e_t+\sqrt kD_t+k^{-1/2}\bigr).
\]
The covariance difference from $C_t$ is at most
$C(e_t+D_t+(t+1)/k)$ entrywise. Indeed a monotone coupling gives
\[
 |\E_QX_sX_r-\E_PX_sX_r|\le d_s(\lambda)+d_r(\lambda),
\]
and the mean correction in the covariance has the same bound up to an
absolute factor. All interpolated covariances consequently remain
uniformly positive definite under the degree condition.

Use the one-face CLT and its Gaussian parameter comparison in the
interpolation. Separating off the change of the root's initial law gives
\begin{align}
 P_{k,\theta}^{u}(\lambda)-P_{k,0}^{u}(\lambda)
 &=\frac{\sqrt k}{2}\sum_{s=0}^t
       d_s(\lambda)\lambda_{s+1}J_s^t(\lambda)+\mathcal E_\lambda,
 \label{eq:bias-path-expansion}\\
 |\mathcal E_\lambda|
 &\le\frac\theta2+
 L_t\sqrt k\left(\sum_{s=0}^td_s(\lambda)\right)
       \bigl(e_t+\sqrt kD_t+k^{-1/2}\bigr),
 \label{eq:bias-path-error}
\end{align}
where $L_t\le\exp\{C2^{t/2}\}$.
The finite prefactor $(k-1)/\sqrt{k-2}$ in the interpolation differs
from $\sqrt k$ by a relative $O(k^{-1})$, which is included in this error.

Multiply \eqref{eq:bias-path-expansion} by $\lambda_{t+1}$ and sum over
root trajectories. The $s=t$ contribution has nonnegative weights:
\[
 \frac{\sqrt k}{2}\sum_{\lambda}d_t(\lambda)J_t^t(\lambda).
\]
By \eqref{eq:response-face-identity}, if
$\eta_t=\sup_u|d_t(u)/b_t-1|$, this lies between
$(1-\eta_t)b_{t+1}$ and $(1+\eta_t)b_{t+1}$.
Thus the previous relative error is not multiplied by a large condition
number.

For $s<t$, use $J_s^t\le(2\pi)^{-1/2}$ and
\eqref{eq:bias-size-bounds}. After division by $b_{t+1}$, the total
absolute contribution is at most
$2^{T+C}\Lambda_T^2/\sqrt k$.
Furthermore, the bootstrap gives $D_t\le12\Lambda_T/k$ and
$\sum_{s\le t}d_s(\lambda)\le3b_t$.
Proposition \ref{prop:unbiased-uniform} bounds $e_t$ by a controlled
constant divided by $\sqrt k$. Summing \eqref{eq:bias-path-error} over
the root trajectories and dividing by $b_{t+1}$ is therefore also at most
a controlled constant divided by $\sqrt k$. The initial-law term has this
bound because $b_t\ge\theta$ and $r_{t+1}^{-1}\le\Lambda_T$.
We have proved
\[
 \eta_{t+1}\le\eta_t+
       \frac{\exp\{C2^{T/2}\}}{\sqrt k}.
\]
Summing over $t$ preserves the controlled-constant form. The degree
condition makes the resulting error less than $1/2$, closing the
bootstrap and proving \eqref{eq:bias-amplification}.
\end{proof}

\section{Saturation and consensus}
\begin{proof}[Proof of Theorem \ref{thm:finite}]
First, the lower bound on $a_T$ and a sufficiently large choice of $A$
ensure that the initialization in \eqref{eq:initial-bias} belongs to $(0,1)$.
At time $T$, \eqref{eq:bias-amplification} gives, uniformly in the parent
trajectory,
\[
 d_T(u)=\frac8{\sqrt k}
   \left(1+O\!\left(\frac{e^{C2^{T/2}}}{\sqrt k}\right)\right).
\]
At earlier times the normalized bias increments are at most a controlled
constant divided by $\sqrt k$, by \eqref{eq:bias-size-bounds}.
Thus the child sum in the exact cavity recursion at horizon $T+1$ has
Gaussian mean
\[
 \mu(\lambda)+8\mathbf e_T+
       O_\infty\!\left(\frac{e^{C2^{T/2}}}{\sqrt k}\right)
\]
and covariance $C_T$ up to an error of the same size.
Here $\mathbf e_T$ denotes the last coordinate unit vector.
Applying the ordinary orthant CLT and Gaussian parameter comparison shows
that the root trajectory through $T+1$, for every imposed parent, differs
in total variation by at most a controlled constant divided by $\sqrt k$
from the cavity process with a field $8$ at its last update only.

Its last spin is
\[
 \sign(H_T+8),\qquad
 H_T=\eta(T)+\sum_{s<T}R(T,s)\tau(s).
\]
Equation \eqref{eq:memory-bounds} gives $\E H_T^2\le4$. Hence
\[
 \E\sign(H_T+8)\ge1-2\Pp(H_T\le-8)\ge1-\frac8{64}=\frac78.
\]
For the degree constant sufficiently large, the approximation error
implies
\begin{equation}\label{eq:rooted-positive}
 \E_{P_{k,\theta}^{u}}\sigma(T+1)\ge\frac12
 \quad\text{for every imposed parent trajectory }u.
\end{equation}

Now return to the full $k$-regular tree and fix a vertex $v$.
In each of the $k$ components obtained by removing $v$, run a lower process
with the parent of that component fixed at $-1$ for all times.
These processes are independent, and their spins at the neighbors of $v$
are simultaneously at most the spins in the original process.
By \eqref{eq:rooted-positive} their means at time $T+1$ are at least $1/2$.
Hoeffding's inequality gives
\[
 \Pp_\theta\{\sigma_v(T+2)=-1\}
 \le\Pp\!\left\{\sum_{i=1}^k\sigma_{i,-}(T+1)\le0\right\}
 \le e^{-k/8}.
\]
This includes possible ties when $k$ is even.
Lemma 1.2 of \cite{KM} says that, for absolute $k_*,\delta_*>0$,
$\E_\theta\sigma_v(t)>1-\delta_*/k$ at some finite time implies
$\theta_*(k)\le\theta$. Our estimate implies this criterion for all
sufficiently large $k$. For each fixed $k$, the selected $T=T(k)$ is
finite, so the use of that criterion entails no uniformity in the
finite-graph limit used in its proof. Finally, the bound on $a_T$ in
\eqref{eq:covariance-lower} proves \eqref{eq:finite-bound}.
\end{proof}

\begin{proof}[Proof of Theorem \ref{thm:asymptotic}]
Let $L=\log k$ and $\ell=\log L$. Choose a fixed, sufficiently large
constant $b$ and set
\[
 T=\left\lfloor\frac2{\log2}(\ell-b)\right\rfloor.
\]
For sufficiently large $k$, this is nonnegative and
\[
 2^{T/2}\le C e^{-b}L.
\]
Choose $b$ so that \eqref{eq:degree-condition} holds. Theorem
\ref{thm:finite} then gives
\begin{align*}
 \log\theta_*(k)
 &\le-\frac{T+1}{2}L+B2^{T/2}\\
 &\le-\frac L{\log2}\ell+O(L).
\end{align*}
When $\theta_*(k)=0$, the same conclusion is read directly as the upper
bound on $\theta_*(k)$ rather than as an assertion about its logarithm.
This proves \eqref{eq:main}.
\end{proof}

\section{Scope of the estimate}
The proof establishes an explicit upper bound for the original dynamics,
not for an independent-neighbor surrogate. It uses monotonicity only in
valid comparisons with imposed parent trajectories. It retains the
response term generated by backtracking. The Gaussian local limit
argument conditions on one coordinate count, and consequently does not
require a lower bound on the probability of every complete trajectory.
The bias comparison uses the exact interpolation
\eqref{eq:product-interpolation}; no neglected higher-order product
expansion is hidden in the critical $k^{-1/2}$ stage.

The estimate does not show that the threshold is positive for large
$k$, does not provide a matching lower bound, and does not identify its
true decay scale. A faster explicit upper bound would follow from stronger
long-time lower bounds on $R(t+1,t)$ and on $\lambda_{\min}(C_T)$.
For example, polynomial-exponential control
$\lambda_{\min}(C_T),r_T\ge e^{-C T^q}$, together with the corresponding
quantitative propagation above, would permit a power-of-$\log k$ time
horizon instead of a $\log\log k$ horizon. Such stronger cavity estimates
are not assumed or proved here.

\begin{thebibliography}{9}
\bibitem{KM}
Y. Kanoria and A. Montanari (2011).
Majority dynamics on trees and the dynamic cavity method.
\emph{The Annals of Applied Probability} \textbf{21}(5), 1694--1748.
DOI: 10.1214/10-AAP729.

\bibitem{Pitt}
L. D. Pitt (1982).
Positively correlated normal variables are associated.
\emph{The Annals of Probability} \textbf{10}(2), 496--499.
DOI: 10.1214/aop/1176993872.

\bibitem{Bentkus}
V. Bentkus (2005).
A Lyapunov-type bound in $\mathbb R^d$.
\emph{Theory of Probability and Its Applications} \textbf{49}(2), 311--323.
DOI: 10.1137/S0040585X97981123.
\end{thebibliography}
\end{document}
